\documentclass{article}
\usepackage{amsmath}
\usepackage{amssymb} % Required for \checkmark

\begin{document}

\[ \text{ex 3 a)} \]
\[ \sum_{k=1}^{6} (7 \cdot 10^{-k}) \rightarrow \sum_{k=x}^{n} c \cdot ak = c \cdot \sum_{k=x}^{n} ak \]
\[ 7 \cdot \sum_{k=1}^{6} (10^{-k}) \rightarrow \sum_{k=0}^{n} (x^k) = \frac{1-x^{n-1}}{1-x} \]
\[ 7 \cdot \left( \frac{1-10^{5-1}}{1-10} \right) \]
\[ \text{On ne part pas de 0 mais de 1} \]
\[ n = 1 \ 2 \ 3 \ 4 \ 5 \ 6 \]
\[ i = 0 \ 1 \ 2 \ 3 \ 4 \ 5 \ 6 \]
\[ D \neq 0 \]
\[ 7 \cdot \left( \frac{1-10^{4}}{-9} \right) \]
\[ 7 \cdot \left( \frac{1-10000}{-9} \right) = 7 \cdot \left( \frac{-9999}{-9} \right) \]
\[ 7 \cdot \frac{9999}{9} \]
\[ 7 \cdot \frac{1111}{1} \]
\[ 7 \cdot 1111 \Rightarrow 7 \cdot 1111 \Rightarrow 7777 \]

\[ \text{ex 3 b)} \]

\[ 5 + 8 + 13 + 17 + \dots + 105 \]

\[ \sum_{k=5}^{105} k + (k-1) \]
\[ \text{Calculation A: } 5+4 = 9 \]
\[ \text{Calculation B: } 6+5 = 11 \]

\[ \sum_{k=5}^{105} \]

\end{document}